\documentclass[11pt]{article}
\usepackage[margin=1.1in]{geometry}
\usepackage{amsmath,amssymb,amsthm,mathtools,booktabs,array,tabularx,longtable}
\usepackage[dvipsnames]{xcolor}
\usepackage[colorlinks=true,linkcolor=MidnightBlue,citecolor=MidnightBlue,urlcolor=MidnightBlue]{hyperref}
\usepackage{fancyhdr}
\pagestyle{fancy}\fancyhf{}\renewcommand{\headrulewidth}{0pt}
\fancyfoot[R]{\thepage}
\fancyfoot[L]{\scriptsize\color{gray}Generated by the Transport Lab from approved parts of its store; not a deposited paper. Papers of record: \href{https://zenodo.org/communities/transport-papers}{Zenodo, Transport Papers}}
\fancypagestyle{plain}{\fancyhf{}\fancyfoot[R]{\thepage}\fancyfoot[L]{\scriptsize\color{gray}Generated by the Transport Lab from approved parts of its store; not a deposited paper. Papers of record: \href{https://zenodo.org/communities/transport-papers}{Zenodo, Transport Papers}}\renewcommand{\headrulewidth}{0pt}}
\newtheorem{theorem}{Theorem}[section]
\newtheorem{proposition}[theorem]{Proposition}
\newtheorem{lemma}[theorem]{Lemma}
\newtheorem{corollary}[theorem]{Corollary}
\newtheorem{conjecture}[theorem]{Conjecture}
\newtheorem{refutation}[theorem]{Refutation}
\newtheorem{prediction}[theorem]{Prediction}
\newtheorem{principle}[theorem]{Principle}
\newtheorem{observation}[theorem]{Observation}
\newtheorem{axiom}[theorem]{Axiom}
\theoremstyle{definition}
\newtheorem{definition}[theorem]{Definition}
\newtheorem{fact}[theorem]{Fact}
\newtheorem{assembled}[theorem]{Assembly}
\newtheorem{postulate}[theorem]{Postulate}
\newtheorem{premise}[theorem]{Premise}
\newtheorem{hypothesis}[theorem]{Hypothesis}
\newtheorem{remark}[theorem]{Remark}
\newtheorem{example}[theorem]{Example}
\theoremstyle{plain}
\setlength{\marginparwidth}{0.8in}
\newcommand{\changedmark}[2]{\marginpar{\raggedright\scriptsize\color{#1}#2}}
\providecommand{\Aa}{}\renewcommand{\Aa}{\mathcal{A}}
\providecommand{\Ad}{}\renewcommand{\Ad}{\widetilde{\mathrm{Ad}}}
\providecommand{\Aev}{}\renewcommand{\Aev}{\mathcal{A}^{+}}
\providecommand{\Aod}{}\renewcommand{\Aod}{\mathcal{A}^{-}}
\providecommand{\B}{}\renewcommand{\B}{\mathrm{B}}
\providecommand{\C}{}\renewcommand{\C}{\mathbb{C}}
\providecommand{\CAR}{}\renewcommand{\CAR}{\mathrm{CAR}}
\providecommand{\CS}{}\renewcommand{\CS}{{\mathbb{S}^{+}}}
\providecommand{\CSm}{}\renewcommand{\CSm}{{\mathbb{S}^{-}}}
\providecommand{\Cl}{}\renewcommand{\Cl}{\mathrm{Cl}}
\providecommand{\Con}{}\renewcommand{\Con}{\mathcal{C}}
\providecommand{\Dl}{}\renewcommand{\Dl}{\Delta}
\providecommand{\EN}{}\renewcommand{\EN}{E_{N}}
\providecommand{\End}{}\renewcommand{\End}{\mathrm{End}}
\providecommand{\F}{}\renewcommand{\F}{\mathbb{F}}
\providecommand{\G}{}\renewcommand{\G}{\Gamma}
\providecommand{\HH}{}\renewcommand{\HH}{\mathbb{H}}
\providecommand{\IS}{}\renewcommand{\IS}{I\sigma_3}
\providecommand{\J}{}\renewcommand{\J}{\mathbf{J}}
\providecommand{\LH}{}\renewcommand{\LH}{L_{H}}
\providecommand{\Mat}{}\renewcommand{\Mat}{\mathrm{M}}
\providecommand{\Mo}{}\renewcommand{\Mo}{M^{\circ}}
\providecommand{\Mp}{}\renewcommand{\Mp}{M^{+}}
\providecommand{\Msun}{}\renewcommand{\Msun}{M_{\odot}}
\providecommand{\Mt}{}\renewcommand{\Mt}{\widetilde{M}}
\providecommand{\Ncone}{}\renewcommand{\Ncone}{\mathcal{N}}
\providecommand{\Neff}{}\renewcommand{\Neff}{N_{\mathrm{eff}}}
\providecommand{\PCS}{}\renewcommand{\PCS}{\mathbb{P}{\mathbb{S}^{+}}}
\providecommand{\PT}{}\renewcommand{\PT}{\mathbb{PT}}
\providecommand{\Pin}{}\renewcommand{\Pin}{\mathrm{Pin}}
\providecommand{\Q}{}\renewcommand{\Q}{\mathbb{Q}}
\providecommand{\R}{}\renewcommand{\R}{\mathbb{R}}
\providecommand{\Ric}{}\renewcommand{\Ric}{\mathrm{Ric}}
\providecommand{\SBH}{}\renewcommand{\SBH}{S_{\mathrm{BH}}}
\providecommand{\Sc}{}\renewcommand{\Sc}{\mathbb{S}}
\providecommand{\Sch}{}\renewcommand{\Sch}{\mathrm{P}}
\providecommand{\Sig}{}\renewcommand{\Sig}{\Sigma}
\providecommand{\Som}{}\renewcommand{\Som}{S_\omega}
\providecommand{\Sp}{}\renewcommand{\Sp}{\mathbb{S}}
\providecommand{\Sph}{}\renewcommand{\Sph}{\mathrm{S}}
\providecommand{\Spi}{}\renewcommand{\Spi}{S_\pi}
\providecommand{\Spin}{}\renewcommand{\Spin}{\mathrm{Spin}}
\providecommand{\Sth}{}\renewcommand{\Sth}{S_{\mathrm{th}}}
\providecommand{\T}{}\renewcommand{\T}{\mathbb{T}}
\providecommand{\TH}{}\renewcommand{\TH}{T_{\mathrm{H}}}
\providecommand{\Tr}{}\renewcommand{\Tr}{\operatorname{Tr}}
\providecommand{\UHF}{}\renewcommand{\UHF}{\mathrm{UHF}}
\providecommand{\Z}{}\renewcommand{\Z}{\mathbb{Z}}
\providecommand{\Zc}{}\renewcommand{\Zc}{\mathrm{Z}}
\providecommand{\cmg}{}\renewcommand{\cmg}{\mathrm{cm^{2}\,g^{-1}}}
\providecommand{\dS}{}\renewcommand{\dS}{\mathrm{dS}}
\providecommand{\dsc}{}\renewcommand{\dsc}{\mathrm{dsc}}
\providecommand{\dtPG}{}\renewcommand{\dtPG}{dt}
\providecommand{\dtS}{}\renewcommand{\dtS}{dt_{\mathrm{S}}}
\providecommand{\ellv}{}\renewcommand{\ellv}{\ell}
\providecommand{\er}{}\renewcommand{\er}{e_r}
\providecommand{\et}{}\renewcommand{\et}{e_t}
\providecommand{\ev}{}\renewcommand{\ev}{\mathrm{ev}}
\providecommand{\gl}{}\renewcommand{\gl}{\mathfrak{gl}}
\providecommand{\hatV}{}\renewcommand{\hatV}{\widehat{V}}
\providecommand{\hb}{}\renewcommand{\hb}{\bar{h}}
\providecommand{\hf}{}\renewcommand{\hf}{\hat{f}}
\providecommand{\id}{}\renewcommand{\id}{\mathrm{id}}
\providecommand{\kpc}{}\renewcommand{\kpc}{\mathrm{kpc}}
\providecommand{\lP}{}\renewcommand{\lP}{\ell_{P}}
\providecommand{\mD}{}\renewcommand{\mD}{\mathbf{m}_{D}}
\providecommand{\mb}{}\renewcommand{\mb}{\bar m}
\providecommand{\mnu}{}\renewcommand{\mnu}{\Sigma m_{\nu}}
\providecommand{\nInf}{}\renewcommand{\nInf}{n_\infty}
\providecommand{\nO}{}\renewcommand{\nO}{n_0}
\providecommand{\ninf}{}\renewcommand{\ninf}{n_{\infty}}
\providecommand{\nz}{}\renewcommand{\nz}{n_{0}}
\providecommand{\proj}{}\renewcommand{\proj}{\operatorname{proj}}
\providecommand{\pv}{}\renewcommand{\pv}{\pi}
\providecommand{\rad}{}\renewcommand{\rad}{\mathrm{rad}}
\providecommand{\sgn}{}\renewcommand{\sgn}{\operatorname{sgn}}
\providecommand{\shA}{}\renewcommand{\shA}{\mathcal{A}}
\providecommand{\shC}{}\renewcommand{\shC}{\mathcal{C}}
\providecommand{\shE}{}\renewcommand{\shE}{\mathcal{E}}
\providecommand{\shJ}{}\renewcommand{\shJ}{\mathcal{J}}
\providecommand{\shK}{}\renewcommand{\shK}{\mathcal{K}}
\providecommand{\shL}{}\renewcommand{\shL}{L}
\providecommand{\shW}{}\renewcommand{\shW}{\mathcal{W}}
\providecommand{\sll}{}\renewcommand{\sll}{\mathfrak{sl}}
\providecommand{\so}{}\renewcommand{\so}{\mathfrak{so}}
\providecommand{\su}{}\renewcommand{\su}{\mathfrak{su}}
\providecommand{\supp}{}\renewcommand{\supp}{\operatorname{supp}}
\providecommand{\tr}{}\renewcommand{\tr}{\operatorname{tr}}
\providecommand{\uu}{}\renewcommand{\uu}{\mathfrak{u}}
\providecommand{\br}{}\renewcommand{\br}[1]{\langle #1\rangle}
\providecommand{\bra}{}\renewcommand{\bra}[1]{\langle #1|}
\providecommand{\file}{}\renewcommand{\file}[1]{\texttt{#1}}
\providecommand{\hp}{}\renewcommand{\hp}[2]{( #1 \,|\, #2 )}
\providecommand{\ip}{}\renewcommand{\ip}[2]{\langle #1, #2\rangle}
\providecommand{\ket}{}\renewcommand{\ket}[1]{|#1\rangle}
\providecommand{\rel}{}\renewcommand{\rel}[2]{D\bigl(#1 \,\|\, #2\bigr)}
\providecommand{\rev}{}\renewcommand{\rev}[1]{\widetilde{#1}}
\providecommand{\slashed}{}\renewcommand{\slashed}[1]{{#1}\mkern-10mu/}
\title{\bfseries The tower and quantum capacity}
\author{Leon Butler\\ \normalsize\href{mailto:lbutler.math@gmail.com}{\texttt{lbutler.math@gmail.com}}\\ \normalsize ORCID: \href{https://orcid.org/0009-0006-4449-5617}{0009-0006-4449-5617}}
\date{Generated by the Transport Lab, 2 November 2026}
\begin{document}
\maketitle
\begin{abstract}
\normalsize
Each time the series repairs a failure of division it adds a pair of new directions, and repeating the step builds a tower of ever larger algebras.

Each floor is the ground algebra with matrix entries, $\Cl(p{+}k, q{+}k) \cong \Mat_{2^k}(\Cl(p,q))$, so nothing essentially new appears. Every floor keeps the ground floor's type, so the clock $d = p - q \bmod 8$ does not advance.

What grows is quantum capacity: each rung adds exactly one fermionic mode. In the infinite limit, for even charge, the tower becomes the algebra of a free fermion field. The walk ends with the ways our spacetime algebra can be extended by a single direction.
\end{abstract}
\newpage
\noindent \textbf{Authorship and disclosure. }The founding idea of the series is the author's: that division by zero in a Clifford algebra should be read not as a prohibition but as a transport into a higher-dimensional algebra, and that the transported division need not be accompanied by an inverse operation. The direction of investigation and the critical review of the results are likewise the author's. The mathematical development and the drafting of the text were carried out in extended collaboration with Claude Fable 5 (Anthropic), an AI system.

\noindent\fbox{\parbox{\dimexpr\linewidth-2\fboxsep-2\fboxrule}{\textbf{This is not a deposited paper.} It was generated by the Transport Lab on 2 November 2026 and assembled from approved parts of its store. It follows results across several papers of the series; each is stated with its proof and credited to the paper that first proved it. The official versions of those papers are on Zenodo, in the Transport Papers community, \url{https://zenodo.org/communities/transport-papers}.}}
\newpage

\section*{Introduction}
The walk starts with the shape of a single floor, then climbs rung by rung, counting what each step adds, and ends at the infinite limit and a side branch. Each result appears after the ones it rests on.

\section{What we assume}

\begin{postulate}[Tameness]\label{c:sheaf.tameness}
This is the standing assumption on how a geometry is allowed to degenerate. The places where it degenerates must form a smooth closed wall, exactly one direction must lose its length at each point of the wall, and these directions must vary smoothly along it. In addition, the length must vanish at a linear rate as the wall is reached, not to higher order.
\end{postulate}

\section{The tower and quantum capacity}

This section climbs the tower one rung at a time.

The tower is built by adding, again and again, a hyperbolic plane: one new positive and one new negative direction.

After $k$ rungs the plane generators give real matrices and a twisted copy of the ground generators commutes with them, so \[ \Cl(p{+}k, q{+}k) \;\cong\; \Cl(p,q) \otimes \Mat\_{2\^{}k}(\R) \] the ground algebra with matrix entries. Climbing enlarges the algebra without changing its character.

\begin{theorem}[Commuting copies]\label{c:tower.commuting-copies}
Adding hyperbolic planes to a Clifford algebra creates nothing essentially new. The enlarged algebra splits into two pieces that commute with each other: a slightly twisted copy of the original algebra, and an algebra of ordinary real square matrices. So each floor of the tower is the ground algebra with matrix entries.
\emph{First proved in} T3, Thm 3.2 (\href{https://doi.org/10.5281/zenodo.22721556}{doi:10.5281/zenodo.22721556}).
\end{theorem}
\noindent\emph{Rests on:} \ref{c:tower.grade-twist-when-it-inner}, \ref{c:tower.volume-elements-planes}, \ref{c:transport.rule-1-morita-trivial}, \ref{c:std-algebra.existence-uniqueness}, \ref{c:std-algebra.blade-basis}.

Because each floor is the ground algebra with matrix entries, its structural features survive the climb.

The charge $d = p - q \bmod 8$, the centre, the division ring and the kinds of obstruction are the same on every floor. For spacetime every floor is \[ \Cl(1{+}k, 3{+}k) \cong \Mat\_{2\^{}{k+1}}(\HH) \] so the clock stays frozen at position $6$.

\begin{theorem}[Rigidity; the clock is frozen]\label{c:tower.rigidity-clock-frozen}
Building the tower by adding pairs of opposite-signed directions never changes the essential type of the algebra. Every floor has the same underlying number system, the same centre and the same kinds of obstruction as the ground floor. For the spacetime algebra, every floor is a matrix algebra over the quaternions, so the eight-step Bott clock does not advance along the tower: it stays where it starts.
\emph{First proved in} T3, Thm 3.6 (\href{https://doi.org/10.5281/zenodo.22721556}{doi:10.5281/zenodo.22721556}).
\end{theorem}
\noindent\emph{Rests on:} \ref{c:tower.commuting-copies}, \ref{c:tower.matrix-amplification-structural-properties}, \ref{c:tower.morita-equivalence-algebra-amplification}, \ref{c:transport.existence-classification-type-c}, \ref{c:std-algebra.sta-matrix}.

If the type never changes, the natural question is what the climb does add.

The null pairs $v_i, w_i$ adjoined at the rungs obey exactly the canonical anticommutation relations of fermionic modes, \[ v\_i v\_j + v\_j v\_i = 0, \qquad w\_i w\_j + w\_j w\_i = 0, \qquad v\_i w\_j + w\_j v\_i = \delta\_{ij} \] as vectors, with no representation chosen. A new rung adds one mode without disturbing the earlier ones.

\begin{proposition}[The ladder pairs satisfy the CAR, as vectors]\label{c:tower.ladder-pairs-satisfy-car-vectors}
Each step up the tower adds a new pair of special zero-length directions to the space. Taken together, these pairs obey exactly the algebraic rules that quantum theory uses for fermionic modes, the kind of states occupied by particles such as electrons, with one mode per step. The rules hold for the directions themselves, so nothing has to be added or chosen to make them work, and a new step never disturbs the modes already present.
\emph{First proved in} T3, Prop 4.1 (\href{https://doi.org/10.5281/zenodo.22721556}{doi:10.5281/zenodo.22721556}).
\end{proposition}
\noindent\emph{Rests on:} \ref{c:std-algebra.vector-squares}, \ref{c:std-algebra.witt-basis-car}.

Fermionic modes need a space of states to act on.

The $v_j$ act by wedging and the $w_j$ by contraction on the exterior algebra $\Lambda^\bullet W_k$ of dimension $2^k$, a Fock space whose vacuum $1$ is killed by every $w_j$. Each rung doubles it, \[ \Lambda\^{}\bullet W\_{k+1} = \Lambda\^{}\bullet W\_k \otimes \Lambda\^{}\bullet \R v\_{k+1} \] one qubit per rung.

\begin{theorem}[Fock module; one qubit per rung]\label{c:tower.fock-module-one-qubit-per-rung}
The planes added in the tower act naturally on a space built like the state space of fermions: an empty vacuum state, with raising operators that fill slots and lowering operators that empty them. Nothing smaller carries this action, and after k planes the space has dimension 2 to the power k. Each new plane doubles it, adding one qubit per rung of the tower.
\emph{First proved in} T3, Thm 4.2 (\href{https://doi.org/10.5281/zenodo.22721556}{doi:10.5281/zenodo.22721556}).
\end{theorem}
\noindent\emph{Rests on:} \ref{c:tower.commuting-copies}, \ref{c:tower.ladder-pairs-satisfy-car-vectors}, \ref{c:std-algebra.existence-uniqueness}, \ref{c:std-algebra.witt-decomposition}.

One can then ask how many fermionic modes an algebra can hold at most, using only its own vectors.

The answer is the Witt index, \[ \min(p,q) \] the smaller of the counts of positive and negative directions. Along the tower it is $\min(p,q) + k$, so each rung adds exactly one mode and no more.

\begin{theorem}[Witt index is quantum capacity]\label{c:tower.witt-index-quantum-capacity}
Inside a Clifford algebra one can ask how many independent fermionic modes can be built directly from vectors. The answer is a classical number attached to the geometry, the Witt index: the largest number of independent zero-length directions that are all orthogonal to one another. For a space with p positive and q negative directions it is the smaller of p and q. Each step up the tower raises this number by exactly one, so each step adds precisely one mode of capacity and no more.
\emph{First proved in} T3, Thm 4.3 (\href{https://doi.org/10.5281/zenodo.22721556}{doi:10.5281/zenodo.22721556}).
\end{theorem}
\noindent\emph{Rests on:} \ref{c:tower.ladder-pairs-satisfy-car-vectors}, \ref{c:std-algebra.vector-squares}, \ref{c:std-algebra.witt-decomposition}.

Having counted rung by rung, the walk now lets the tower run forever.

When $p - q$ is even the algebraic union has a clean form, \[ \varinjlim T\_k \;\cong\; \Cl(p,q) \otimes \Mat\_{2\^{}\infty} \] the ground algebra with entries in an infinite matrix algebra built by repeated doubling. Equivalently, it is the Clifford algebra of $\R^{p,q}$ with all the planes added at once.

\begin{theorem}[The algebraic limit, even charge]\label{c:tower.algebraic-limit-even-charge}
Keep adding planes to the tower forever and take the union, purely algebraically. When the numbers of positive and negative directions differ by an even amount, the result has a clean description: the original algebra with entries in an infinite matrix algebra built by repeated doubling. Equivalently, it is the Clifford algebra of the infinite-dimensional space obtained by adding all the planes at once.
\emph{First proved in} T3, Thm 5.3 (\href{https://doi.org/10.5281/zenodo.22721556}{doi:10.5281/zenodo.22721556}).
\end{theorem}
\noindent\emph{Rests on:} \ref{c:tower.commuting-copies}, \ref{c:tower.compatibility-connecting-maps-transported}.

The algebraic union still has to be completed, with complex scalars and a norm.

For even charge the completed limit is \[ \Cl(p,q)\_\C \,\otimes\, \CAR \] the ground algebra with the CAR algebra, the standard algebra of a free fermion field with countably many modes. The null vectors added at each rung act directly as that field's creation and annihilation operators, with $w_j = v_j^*$.

\begin{theorem}[The CAR limit, even charge]\label{c:tower.car-limit-even-charge}
Repeatedly applying the transport step builds an infinite tower of algebras. When the signature difference is even and the tower is completed, it becomes the original algebra combined with the standard algebra of a free field of fermions with countably many modes. Each rung of the tower contributes exactly one fermionic mode, and the null vectors added at that rung act directly as its creation and annihilation operators.
\emph{First proved in} T3, Thm 5.6 (\href{https://doi.org/10.5281/zenodo.22721556}{doi:10.5281/zenodo.22721556}).
\end{theorem}
\noindent\emph{Rests on:} \ref{c:tower.algebraic-limit-even-charge}, \ref{c:tower.fock-module-one-qubit-per-rung}, \ref{c:tower.glimms-classification-car-uniqueness}, \ref{c:tower.ladder-pairs-satisfy-car-vectors}, \ref{c:tower.structure-tower-c-norms}.

These results answer the question that started the tower: how many times must the completion step be repeated?

Under the tameness premise, for a compact wall with a trivial radical line bundle, one ascent completes the physical side, and further rungs arise only beyond it, one per shadow wall.

The repetitions never change the algebra's type, and in the even-charge limit they add the modes of a free fermion field.

\begin{corollary}[Answers to the iteration question]\label{c:tower.answers-iteration-question}
The question was how many times the series' completion step has to be repeated. Under the tameness premise, and for a compact wall with a trivial radical line, on the physical side of a degeneration the answer is once; further steps are demanded only beyond it, one for each successive hidden wall. Repeating the step never changes the algebra's essential type, and in the infinite limit, for even charge, it simply adds the modes of a free fermion field, one qubit per step.
\emph{First proved in} T3, Cor 6.2 (\href{https://doi.org/10.5281/zenodo.22721556}{doi:10.5281/zenodo.22721556}).
\end{corollary}
\noindent\emph{Rests on:} \ref{c:tower.wall-package}, \ref{c:tower.rigidity-clock-frozen}, \ref{c:tower.algebraic-limit-even-charge}, \ref{c:tower.car-limit-even-charge}, \ref{c:sheaf.tameness}.

The unit the tower adds turns up elsewhere in the series under other names.

A ladder pair $(v_i, w_i)$, a fermionic mode and a unit of Witt-index capacity are one object: a single qubit, of maximal entropy $\ln 2$ nats. The conformal algebra $\Cl(2,4)$ carries exactly two of them, whether one counts pairs or modes.

\begin{proposition}[Pair, mode, qubit: one object]\label{c:scrambling.pair-mode-qubit-one-object}
Three things in the series that sound different, a pair of null directions in the tower, a fermionic mode in the singularity analysis and a unit of quantum capacity, are the same thing. Each is one fermion mode with two states, empty or filled, which is one qubit and can hold at most $\ln 2$ nats of entropy. The conformal algebra has room for exactly two of them, and counting pairs or counting modes gives the same number.
\emph{First proved in} T8, Prop 4.11 (\href{https://doi.org/10.5281/zenodo.22722176}{doi:10.5281/zenodo.22722176}).
\end{proposition}
\noindent\emph{Rests on:} \ref{c:tower.fock-module-one-qubit-per-rung}, \ref{c:tower.ladder-pairs-satisfy-car-vectors}, \ref{c:tower.witt-index-quantum-capacity}, \ref{c:std-algebra.witt-basis-car}.

So far each step added two directions at once, and adding a single direction instead leads off the main tower.

Our spacetime algebra $\Cl(1,3)$ and its mirror $\Cl(3,1)$ can each be extended by one time or one space direction, giving $\Cl(1,4)$, $\Cl(2,3)$, $\Cl(3,2)$ and $\Cl(4,1)$. A space direction doubles, as in $\Cl(1,4) = \Cl(1,3) \oplus \Cl(1,3)$, a time direction complexifies, and the four land on the residues $5$, $7$, $1$ and $3$.

\begin{corollary}[The four extensions of the oriented pair]\label{c:odd-extensions.four-extensions-oriented-pair}
Our spacetime algebra and its mirror image can each be extended by one extra direction, of time type or of space type, giving four algebras. Adding a space direction produces two identical copies of the original algebra side by side, while adding a time direction produces the original algebra with complex numbers allowed. Each of the four lands on a different odd value of the series' mod-eight charge.
\emph{First proved in} T14a, Cor 2.4 (\href{https://doi.org/10.5281/zenodo.23232104}{doi:10.5281/zenodo.23232104}).
\end{corollary}
\noindent\emph{Rests on:} \ref{c:odd-extensions.doubling-or-complexification}, \ref{c:odd-extensions.imported-representatives-dimension-nearest-four-types}, \ref{c:odd-extensions.volume-element-odd-dimension}, \ref{c:std-algebra.sta-matrix}.

\section{Results used from elsewhere}
The proofs above use the following results of the series, stated briefly; each is proved in the paper named.

\begin{fact}[Doubling or complexification]\label{c:odd-extensions.doubling-or-complexification}
An odd-dimensional $\Cl(p,q)$ is two copies of its even subalgebra when $\omega^{2} = +1$ and that subalgebra complexified when $\omega^{2} = -1$, and the even subalgebra is a Clifford algebra of one dimension fewer.
\emph{Source:} T14a, Thm 2.3 (\href{https://doi.org/10.5281/zenodo.23232104}{doi:10.5281/zenodo.23232104}).
\end{fact}

\begin{fact}[Imported: the representatives at the dimension nearest four, and the types]\label{c:odd-extensions.imported-representatives-dimension-nearest-four-types}
Imported: the nearest-four representatives of the eight residues, with $r = 0, 4$ sterile and the other six admitting matter, and spinor type a function of the residue, so that $\Cl(1,3) \cong \Mat(2,\HH)$ and $\Cl(3,1) \cong \Mat(4,\R)$ differ.
\emph{Source:} T14a, Fact 1.2 (\href{https://doi.org/10.5281/zenodo.23232104}{doi:10.5281/zenodo.23232104}).
\end{fact}

\begin{fact}[The volume element in odd dimension]\label{c:odd-extensions.volume-element-odd-dimension}
In odd dimension $n$ the volume element of $\Cl(p,q)$ is central with $\omega^{2} = (-1)^{n(n-1)/2 + q}$, which is $(-1)^{q}$ for $n = 5$.
\emph{Source:} T14a, Lem 2.2 (\href{https://doi.org/10.5281/zenodo.23232104}{doi:10.5281/zenodo.23232104}).
\end{fact}

\begin{fact}[The blade basis, degenerate forms included]\label{c:std-algebra.blade-basis}
For any basis of $V$ the $2^{n}$ increasing products $e_I$ form a basis of $\Cl(V, Q)$, degenerate forms included; so $V \hookrightarrow \Cl(V, Q)$.
\emph{Source:} the algebra core of the Transport Lab, proved there with an exact check; standard, see C. Chevalley, The Algebraic Theory of Spinors, Columbia University Press (1954), Ch. II (the Clifford algebra of a quadratic space).
\end{fact}

\begin{fact}[The Clifford algebra exists and is unique]\label{c:std-algebra.existence-uniqueness}
$\Cl(V, Q) = T(V)/\langle v \otimes v - Q(v) \rangle$ exists and is unique up to unique isomorphism; the wedge-plus-contraction operators on $\Lambda V$ give a concrete representation.
\emph{Source:} the algebra core of the Transport Lab, proved there with an exact check; standard, see C. Chevalley, The Algebraic Theory of Spinors, Columbia University Press (1954), Ch. II (the Clifford algebra of a quadratic space).
\end{fact}

\begin{fact}[The spacetime algebra as a matrix algebra, and the spacetime split]\label{c:std-algebra.sta-matrix}
$\Cl(1,3) \cong \Mat_2(\HH)$, the Dirac matrices represent $\C \otimes \Cl(1,3) \cong \Mat_4(\C)$, and $\sigma_k = \gamma_k\gamma_0$ give $\Cl^{+}(1,3) \cong \Cl(3,0)$ with $\sigma_1\sigma_2\sigma_3 = I$.
\emph{Source:} the algebra core of the Transport Lab, proved there with an exact check; standard, see C. Doran and A. Lasenby, Geometric Algebra for Physicists, Cambridge University Press (2003), Ch. 5 (relativity and spacetime).
\end{fact}

\begin{fact}[The product of two vectors]\label{c:std-algebra.vector-squares}
$v^{2} = Q(v)$ and $uv + vu = 2B(u, v)$, so $uv = B(u, v) + \tfrac12(uv - vu)$.
\emph{Source:} the algebra core of the Transport Lab, proved there with an exact check; standard, see C. Doran and A. Lasenby, Geometric Algebra for Physicists, Cambridge University Press (2003), Ch. 4 (foundations of geometric algebra).
\end{fact}

\begin{fact}[Null pairs satisfy the canonical anticommutation relations]\label{c:std-algebra.witt-basis-car}
For a Witt basis ($B(v_i,v_j) = B(w_i,w_j) = 0$, $B(v_i,w_j) = \tfrac12\delta_{ij}$), $\{v_i,v_j\} = \{w_i,w_j\} = 0$ and $\{v_i,w_j\} = \delta_{ij}$ in $\Cl(V,Q)$.
\emph{Source:} the algebra core of the Transport Lab, proved there with an exact check; standard, see C. Doran and A. Lasenby, Geometric Algebra for Physicists, Cambridge University Press (2003), Ch. 9 (multiparticle states).
\end{fact}

\begin{fact}[Hyperbolic pairs and the Witt decomposition]\label{c:std-algebra.witt-decomposition}
A null vector in a nondegenerate plane has a unique hyperbolic partner, and $V = H_1 \perp \dots \perp H_{\min(p,q)} \perp A \perp \operatorname{rad} V$ with $A$ definite; maximal totally null subspaces have dimension $\min(p,q) + r$.
\emph{Source:} the algebra core of the Transport Lab, proved there with an exact check; standard, see T. Y. Lam, Introduction to Quadratic Forms over Fields, Graduate Studies in Mathematics 67, American Mathematical Society (2005), Ch. I (hyperbolic planes, the Witt decomposition, Witt's theorems).
\end{fact}

\begin{fact}[Compatibility: the connecting maps, transported]\label{c:tower.compatibility-connecting-maps-transported}
For even charge the transported tower embedding is block doubling followed by conjugation by the invertible corrector $V_k = \operatorname{diag}(u_A \otimes 1_{2^k},\ 1_A \otimes \overline\Omega_k)$, which splits the correction across the two tensor legs.
\emph{Source:} T3, Lem 5.2 (\href{https://doi.org/10.5281/zenodo.22721556}{doi:10.5281/zenodo.22721556}).
\end{fact}

\begin{fact}[Glimm's classification, and CAR uniqueness]\label{c:tower.glimms-classification-car-uniqueness}
The completion of a unital chain of matrix algebras $\Mat_{2^k}(\C)$ is the unique $2^\infty$ UHF algebra, and the CAR algebra over $\ell^2$ is canonically determined by its relations and isomorphic to it.
\emph{Source:} J. Glimm, On a certain class of operator algebras, Trans. Amer. Math. Soc. 95 (1960), 318--340.; O. Bratteli and D. W. Robinson, Operator Algebras and Quantum Statistical Mechanics 2, 2nd ed., Springer, Berlin (1997)., Thm.~5.2.5.
\end{fact}

\begin{fact}[The grade twist, and when it is inner]\label{c:tower.grade-twist-when-it-inner}
The tower embedding is a diagonal twisted by the grade involution, which is inner, implemented by the ground volume element, when $p - q$ is even, and outer when $p - q$ is odd.
\emph{Source:} T3, Lem 5.1 (\href{https://doi.org/10.5281/zenodo.22721556}{doi:10.5281/zenodo.22721556}).
\end{fact}

\begin{fact}[Matrix amplification: structural properties]\label{c:tower.matrix-amplification-structural-properties}
Matrix amplification $A \mapsto \Mat_m(A)$ preserves the centre and simplicity, regroups nested matrix algebras, and distributes over direct sums.
\emph{Source:} T3, Lem 3.3 (\href{https://doi.org/10.5281/zenodo.22721556}{doi:10.5281/zenodo.22721556}).
\end{fact}

\begin{fact}[Morita equivalence of an algebra and its amplification]\label{c:tower.morita-equivalence-algebra-amplification}
Every unital ring $A$ is Morita equivalent to the matrix ring $\Mat_m(A)$, via the column module $A^m$.
\emph{Source:} T. Y. Lam, Lectures on Modules and Rings, Graduate Texts in Mathematics 189, Springer, New York (1999)., \S\S 17--18.
\end{fact}

\begin{fact}[The *-structure of the tower, and its C*-norms]\label{c:tower.structure-tower-c-norms}
Each complexified floor of the tower has a unique frame-adapted *-structure respected by all the tower's maps and correctors, and for $p - q$ even each floor is a full complex matrix algebra, so the floors and their union carry unique C*-norms.
\emph{Source:} T3, Lem 5.4 (\href{https://doi.org/10.5281/zenodo.22721556}{doi:10.5281/zenodo.22721556}).
\end{fact}

\begin{fact}[Volume elements of the planes]\label{c:tower.volume-elements-planes}
The plane volume elements $\omega_j = e_j f_j$ square to one, anticommute with their own plane, commute with all other generators and with each other, and their product $\Omega_k$ squares to one, commutes with the ground generators and anticommutes with every plane generator.
\emph{Source:} T3, Lem 3.1 (\href{https://doi.org/10.5281/zenodo.22721556}{doi:10.5281/zenodo.22721556}).
\end{fact}

\begin{fact}[The wall package]\label{c:tower.wall-package}
Under tameness and a trivial radical line bundle the one-sided completion is nondegenerate through the wall and relocates the degeneracy to a shadow wall; for a compact wall and good couplings this iterates, adding one hyperbolic plane per wall.
\emph{Source:} T3, Fact 6.1 (\href{https://doi.org/10.5281/zenodo.22721556}{doi:10.5281/zenodo.22721556}).
\end{fact}

\begin{fact}[Existence and classification of type C]\label{c:transport.existence-classification-type-c}
$\Cl(p,q)$, $n \geq 1$, has a nontrivial central idempotent exactly when $p - q \equiv 1 \pmod 4$, and then the only ones are $\varepsilon_\pm = \tfrac12(1 \pm \omega)$, which split the algebra into two simple ideals and are exchanged by the grade involution but by no inner automorphism.
\emph{Source:} T1, Thm 3.6 (\href{https://doi.org/10.5281/zenodo.21743863}{doi:10.5281/zenodo.21743863}).
\end{fact}

\begin{fact}[Rule 1 is Morita-trivial]\label{c:transport.rule-1-morita-trivial}
Rule 1 gives $\Cl(p+1, q+1) \cong \Mat_2(\Cl(p,q))$ with the adjoined ladder pair as off-diagonal matrix units, so it preserves the Morita class, the Wedderburn type and the clock residue $p - q \pmod 8$.
\emph{Source:} T1, Thm 7.1 (\href{https://doi.org/10.5281/zenodo.21743863}{doi:10.5281/zenodo.21743863}).
\end{fact}

\section*{Conclusion}
The tower never changes an algebra's type, yet each rung contributes exactly one fermionic mode, and for even charge the infinite climb yields a free fermion field.

The account of when further rungs are needed rests on the tameness assumption listed under what we assume. That premise is the only assumption this trail lists.

\end{document}